Zur Umsetzung unserer Lösungsidee benötigten wir neben den Gang- und Raumszenen zusätzlich ein Pfeilobjekt und ein Türobjekt, welche den Wechsel zwischen verschiedenen Szenen möglich machen.

\subsection{Das Pfeilobjekt}
\begin{sloppypar}
Um den Szenenwechsel per Pfeil zu ermöglichen, wurde eine Klasse mit dem Namen \texttt{PfeilObject} angelegt. Dabei müssen individuell für jede Szene die Parameter Originalszene (Szene, in der man sich gerade befindet), Zielszene (Szene, in die der Pfeil führt) und die Richtung des Pfeils (rechts, links, oben, unten) veränderbar sein; das sind also die formalen Parameter des Konstruktors der Klasse.

Innerhalb der Methode wurde mittels eines Switch-Case-Blocks für die verschiedenen Richtungen der Pfeile passende Koordinaten für die jeweilige Position eines Pfeils festgelegt. Für die verschiedenen Richtungen der Pfeile (rechts, links, oben, unten) haben wir vier verschiedene Bilddateien verwendet. \\
Damit am Ende bei der Richtungsangabe das richtige Bild verwendet wird, haben wir eine \texttt{getSprite}-Methode erstellt, in der je nach Angabe einer der vier Richtungen der Sprite des richtigen Bildes zurückgegeben wird. Diese wird im Konstruktor aufgerufen. \\
Am Ende der Klasse wird in einer \texttt{update}-Methode die alte Szene mit der durch das \texttt{PfeilObject} neu festgelegten Szene im Szenenstapel ersetzt.
\end{sloppypar}

\begin{lstlisting}[caption=PfeilObject, label=türobject, captionpos=t, language=Java]
public class PfeilObject extends BaseObject{
		
	public static final int RECHTS = 1, LINKS = 2, 	OBEN = 3,
	UNTEN = 4;
		
	/* Initialisieren eines Ziels (Wohin soll der Pfeil 
	fuehren?) */
	public Supplier<BaseScene> target;
	//Initialisieren eines Ursprungs
	public BaseScene origin;
		
	public PfeilObject(BaseScene origin, Supplier<BaseScene> 
	target, int richtung){
		/* Festlegen, welches Bild genutzt wird (siehe 
		weiter unten) */
		super(getSprite(richtung));
			
		/* Festlegen der Koordinaten des Objektes abhaengig 
		davon,	welche Richtung angegeben wird */
		switch(richtung){
			case RECHTS:
				x = 880;
				y = (HEIGHT/2) - 60;
				break;
			case LINKS:
				x = 9;
				y = (HEIGHT/2) - 60;
				break;
			case OBEN:
				x = (WIDTH/2) - 40;
				y = 4;
				break;
			case UNTEN:
				x = (WIDTH/2) - 40;
				y = 452;
				break;
		}
			
		this.target = target;
		this.origin = origin;
			
			
	}
		
		
	/* Funktion fuer die Bilder, je nachdem, welche Richtung 
	beim PfeilObject angegeben wird, wird ein anderes Bild 
	fuer das PfeilObject genutzt */
	public static Image getSprite(int richtung) {
		switch(richtung){
			case RECHTS:
				return Images.get("/assets/background/PfeilRechts.png");
			case LINKS:
				return Images.get("/assets/background/PfeilLinks.png");
			case OBEN:
				return Images.get("/assets/background/PfeilOben.png");
			case UNTEN:
				return Images.get("/assets/background/PfeilUnten.png");
		}
		return null;
	}
		
	@Override
	public void update(BaseScene scene) {
		if(isClicked()){
			BaseScene neu = target.get();
			SceneStack.INSTANCE.replace(origin, neu);
		}
	}
}
\end{lstlisting}

\subsection{Das Türobjekt}
Für den Wechsel von einer Gang- in eine Raumszene gibt es keinen Pfeil, sondern ein Türobjekt, das sogenannte \texttt{InvisibleDoorObject}. Von einem Raum aus kommt man mit einem Pfeil wieder in den Gang zurück. \\
Dieses Türobjekt ist ähnlich wie das Pfeilobjekt gestaltet, mit dem Unterschied, dass das Objekt keine Richtungsangabe, sondern eine Position benötigt und kein Bild verwendet. Stattdessen besitzt das Objekt nur eine Größe, die den Maßen der Türen in den Hintergründen entspricht, sodass das Objekt \glqq unsichtbar\grqq\ ist und dann auf die Koordinaten der Tür im Hintergrundbild gelegt werden kann. \\
Man klickt im Endeffekt also nicht tatsächlich auf die Tür, sondern auf ein unsichtbares Objekt, welches über die Tür gelegt wurde, wobei für den Spieler dieser Unterschied nicht erkennbar ist. \\
In allen anderen Bereichen funktioniert das Türobjekt genauso wie das Pfeilobjekt, es werden Original- und Zielszene festgelegt, welche am Ende im Szenenstapel entfernt (Originalszene) bzw. hinzugefügt (Zielszene) werden.

\begin{lstlisting}[caption=InvisibleDoorObject, label=lst:invisibledoorobject, captionpos=t, language=Java]
public class InvisibleDoorObject extends BaseObject{
	//Initialisieren eines Ziels (Wohin soll die Tuer fuehren?)
	public Supplier<BaseScene> target;
	//Initialisieren eines Ursprungs
	public BaseScene origin;
		
	public InvisibleDoorObject(int x, int y, BaseScene origin, 
	Supplier<BaseScene> target){
		//Einfuegen des Objektes
		//kein Bild (Objekt ist unsichtbar)
		/* x und y beschreiben die Position des Objektes, 
		diese wird in den jeweiligen Szenen festgelegt */
		/* Breite und Hoehe entsprechen den Werten der Tuer 
		in den Bildern */
		super(null, x, y, 140, 240);
		this.target = target;
		this.origin = origin;
	}
		
	@Override
	public void update(BaseScene scene) {
		if(isClicked())
		SceneStack.INSTANCE.replace(origin, target.get());
	}
}
\end{lstlisting}

\subsection{Gang- und Raumszenen}
Für jede Gang-/Raumszene wurde eine eigene Klasse erstellt, welche allgemein mit \glqq OrtScene\grqq\ benannt wurden. In diesen Klassen, welche die \texttt{BaseScene} erweitern, wird zunächst auf das jeweilige Bild zugegriffen. Danach wird die Pfeilobjekt-Klasse verwendet, um Pfeilobjekte einzufügen, welche in andere Szenen führen. Zum Schluss wird die Spielfigur mit in jeder Szene individueller Position zur Szene hinzugefügt.

\begin{lstlisting}[caption=Mathegangszene, label=lst:mahtegangszene, captionpos=t, language=Java]
public class MathegangScene extends BaseScene{
	public MathegangScene(){
		//Einfügen des Hintergrunds
		objects.add(new  BackdropObject(Images.get(
			"/assets/background/Mathegang.png")));
		
		/* Einfügen des unsichtbaren Tür-Objektes auf 
		Position der Tür */
		objects.add(new InvisibleDoorObject(
			MainWindow.PIXEL_SCALE * 177, MainWindow.PIXEL_SCALE * 45, this, MatheraumScene::new));
		
		/* Pfeilobjekte für den Wechsel in nebenliegende Szenen */
		objects.add(new PfeilObject(this, 
			TreppeRechts2Scene::new, PfeilObject.RECHTS));
		objects.add(new PfeilObject(this, 
			ErstesOGScene::new, PfeilObject.LINKS));
		
		//Einfügen der Spieler-Figur
		this.objects.add(PlayerObject.INSTANCE);
		PlayerObject.INSTANCE.x = PlayerObject.GetRandomX();
		PlayerObject.INSTANCE.y = 295;
	}
}
\end{lstlisting}

\subsubsection{Position der Spielfigur}
In den meisten Szenen wird die x-Koordinate mithilfe der Funktion \texttt{GetRandomX} bestimmt, welche in der Klasse \texttt{PlayerObject} festgelegt ist. 

\begin{lstlisting}[caption=GetRandonX-Methode, label=lst:getrandomx, captionpos=t, language=Java]
/* Funktion zum Festlegen einer zufaelligen x-Koordinate fuer die 
Spielfigur */
public static int GetRandomX(){
	Random r = new Random();
	int min = 35;
	int max = 790;
	int position = r.nextInt((max - min) + 1) + min;
	return position;
}
\end{lstlisting}

In anderen Szenen (wie zum Beispiel die Foyerszene) wurde auch die x-Koordinate fest definiert, da es sonst Komplikationen mit dem Hintergrunddesign gegeben hätte. So soll die Spielfigur zum Beispiel nicht auf einem Blumentopf stehen.
